\hnsection{连续群与无穷小变换}

1873 年秋，Lie 突然重新开始研究变换群，并取得决定性成果。根据他 1873--1874 年写给 A. Mayer 的几封信（{[}3{]}，第 5 卷，第 584--608 页），在能够追踪其思路的范围内，他从作用于 n 个变量的“连续群”变换出发

\[
x _ {i} ^ {\prime} = f _ {i} (x _ {1}, \dots , x _ {n}, a _ {1}, \dots , a _ {r}) \quad (1 \leqslant i \leqslant n)\tag{4}
\]

它有效地†依赖于 r 个参数 \(a_{1},\ldots,a_{r}\)；他注意到，若变换 (4) 在参数取 \(a_{1}^{0},\ldots,a_{r}^{0}\) 时是恒等变换‡，则 \(x_{i}\) 的一阶 Taylor 展开：

\[
\begin{array}{r l} f _ {i} (x _ {1}, \dots , x _ {n}, a _ {1} ^ {0} + z _ {1}, \dots , a _ {r} ^ {0} + z _ {r}) & \\ = x _ {i} + \sum_ {k = 1} ^ {r} z _ {k} X _ {k i} (x _ {1}, \dots , x _ {n}) + \dots & (1 \leqslant i \leqslant n) \end{array} \tag {5}
\]

给出一个关于 r 个参数 \(z_{j}\) 线性依赖的“通用”无穷小变换

\[
d x _ {i} = \left(\sum_ {k = 1} ^ {r} z _ {k} \mathrm{X} _ {k i} (x _ {1}, \dots , x _ {n})\right) d t \quad (1 \leqslant i \leqslant n).\tag{6}
\]

仿照他与 Klein 合写的论文，Lie 对微分系统作了积分

\[
\frac {d \xi_ {1}}{\sum_ {k} z _ {k} X _ {k 1} (\xi_ {1} , \dots , \xi_ {n})} = \dots = \frac {d \xi_ {n}}{\sum_ {k} z _ {k} X _ {k n} (\xi_ {1} , \dots , \xi_ {n})} = d t,\tag{7}
\]

这使他对每个点 \((z_{1},\ldots,z_{r})\) 得到一个单参数群

\[
t \mapsto x _ {i} ^ {\prime} = g _ {i} (x _ {1}, \dots , x _ {n}, z _ {1}, \dots , z _ {r}, t) \quad (1 \leqslant i \leqslant n)\tag{8}
\]

从而 \(g_{i}(x_{1},\ldots,x_{n},z_{1},\ldots,z_{r},0)=x_{i}\) 对所有 i 成立。他用一种巧妙的方法，利用变换 (4) 在复合下构成稳定集合这一事实，证明单参数群 (8) 是给定群的子群 {[}3 d{]}。整个理论的关键新思想，是把函数 (4) 的 Taylor 展开取到二阶。他的论证过程相当混乱而带有启发式色彩（{[}3 d{]} 和 {[}3{]}，第 5 卷，第 600--601 页）；可以叙述如下。对于足够小的 \(z_{r}\)，令 (8) 中 t=1；于是得到该群变换的新参数 \(z_{1},\ldots,z_{r}\)（这实际上是“典范参数”一词首次出现）。然后根据 (7)，

\[
\frac {\partial g _ {i}}{\partial t} = \sum_ {k} z _ {k} \mathrm{X} _ {k i} (x _ {1} ^ {\prime}, \dots , x _ {n} ^ {\prime})
\]

由此

\[
\begin{array}{r l} \frac {\partial^ {2} g _ {i}}{\partial t ^ {2}} & = \sum_ {k, j} z _ {k} \frac {\partial \mathrm{X} _ {k i}}{\partial x _ {j}} (x _ {1} ^ {\prime}, \ldots , x _ {n} ^ {\prime}) \frac {\partial x _ {j} ^ {\prime}}{\partial t} \\ & = \sum_ {k, j} z _ {k} \frac {\partial \mathrm{X} _ {k i}}{\partial x _ {j}} (x _ {1} ^ {\prime}, \ldots , x _ {n} ^ {\prime}) \left(\sum_ {n} z _ {h} \mathrm{X} _ {h j} (x _ {1} ^ {\prime}, \ldots , x _ {n} ^ {\prime})\right) \end{array}
\]

从而得到

\[
\begin{array}{r} x _ {i} ^ {\prime} = x _ {i} + \Big (\sum_ {k} z _ {k} \mathrm{X} _ {k i} (x _ {1}, \ldots , x _ {n}) \Big) t \\ + \frac {1}{2} \Big (\sum_ {k, n, j} z _ {k} z _ {h} \frac {\partial \mathrm{X} _ {k i}}{\partial x _ {j}} (x _ {1}, \ldots , x _ {n}) \mathrm{X} _ {h j} (x _ {1}, \ldots , x _ {n}) \Big) t ^ {2} + \dots , \end{array}
\]

由此，当 \(t = 1\) 时，关于参数 \(z_{j}\) 的 Taylor 展开为

\[
x _ {i} ^ {\prime} = x _ {i} + \left(\sum_ {j} z _ {k} \mathrm{X} _ {k i}\right) + \frac {1}{2} \left(\sum_ {k, h, j} z _ {k} z _ {h} \mathrm{X} _ {h j} \frac {\partial \mathrm{X} _ {k i}}{\partial x _ {j}}\right) + \dots (1 \leqslant i \leqslant n).\tag{9}
\]

将这些关系简写为 \(x' = \mathbf{G}(x, z)\)，用来联系下列向量

\[
x = (x _ {1}, \dots , x _ {n}), \quad x ^ {\prime} = (x _ {1} ^ {\prime}, \dots , x _ {n} ^ {\prime}), \quad z = (z _ {1}, \dots , z _ {r});
\]

这些变换集合关于复合的基本稳定性可写成

\[
\mathrm{G} (\mathrm{G} (x, u), v) = \mathrm{G} (x, \mathrm{H} (u, v))\tag{10}
\]

其中 \(\mathbf{H} = (\mathbf{H}_1, \ldots, \mathbf{H}_r)\) 与 \(x\) 无关；显然 \(\mathrm{H}(u, 0) = u\) 且 \(\mathrm{H}(0, v) = v\)，从而展开式为

\[
\mathrm{H} _ {i} (u, v) = u _ {i} + v _ {i} + \frac {1}{2} \sum_ {h, k} c _ {i k h} u _ {h} v _ {k} + \dots ,\tag{11}
\]

其中省略的项不关于 u 或 v 线性。将 (10) 与 (9)、(11) 结合变换，然后比较两边关于 \(u_{h}v_{k}\) 的项，Lie 得到关系

\[
\sum_ {j = 1} ^ {n} \left(\mathrm{X} _ {h j} \frac {\partial \mathrm{X} _ {k i}}{\partial x _ {j}} - \mathrm{X} _ {k j} \frac {\partial \mathrm{X} _ {h i}}{\partial x _ {j}}\right) = \sum_ {l = 1} ^ {r} c _ {l h k} \mathrm{X} _ {l i} (1 \leqslant h, k \leqslant r, 1 \leqslant i \leqslant n). \tag{12}
\]

他在偏微分方程理论方面的经验使他把这些条件写成更简单的形式：按照 (1) 的模式，对 (6) 中令 \(z_{k} = 1\)、\(z_{h} = 0\)（\(h \neq k\)）所得到的每个无穷小变换，赋予微分算子

\[
\mathrm{A} _ {k} (f) = \sum_ {i = 1} ^ {n} \mathrm{X} _ {k i} \frac {\partial f}{\partial x _ {i}},\tag{13}
\]

并将条件 (12) 改写为

\[
[ \mathrm{A} _ {h}, \mathrm{A} _ {k} ] = \sum_ {l} c _ {l h k} \mathrm{A} _ {l},\tag{14}
\]

这成为他理论的基石。在此之前，他不加区分地使用“无限小变换”和“无穷小变换”这两个术语（例如 {[}3 e{]}）；关系 (14) 的简洁性使他把无穷小变换 \(dx_{i} = X_{ki}dt\)（\(1 \leqslant i \leqslant n\)）的算子 (13) 称为“符号” {[}3 d{]}，很快又把算子 (13) 本身称为“无穷小变换”（{[}3 d{]} 和 {[}3{]}，第 5 卷，第 589 页）。

随后，他意识到“连续群”理论与自己早先关于接触变换和偏微分方程的研究之间存在密切联系。这种汇合使他充满热情：1874 年他在写给 Mayer 的信中说：“我早期的工作仿佛早已在那里，等待着创立变换群的新理论”（{[}3{]}，第 5 卷，第 586 页）。

在接下来的几年里，Lie 继续研究变换群。除下文（§ III）概述的一般定理外，他还得到了一些更特殊的结果：确定直线和平面的变换群、射影群中的低余维子群、至多含 6 个参数的群，等等。他并没有长期放弃微分方程。事实上，对他而言，变换群理论似乎是积分微分方程的一种工具，其中变换群所起的作用类似于代数方程的 Galois 群。†

我们还要指出，这项研究也促使他引入某些具有无穷多个参数的变换集合，他称之为“无限连续群” \(^{‡}\)；他把上面 (4) 型、具有有限个参数的变换群称为“有限连续群”。

